% !TEX root = ../manuscript.tex

\section{Introduction}

Faults are ubiquitous in sedimentary basins and may either impede or promote
subsurface fluid migration. In geologic carbon sequestration, this dual role
creates uncertainty in whether injected \coTwo{} remains in the intended
storage interval, enters a fault, migrates toward or into the top seal, or
moves upward along the fault. Quantifying these possibilities is important for
site selection, injection management, and monitoring design.

The effective flow behavior of a fault emerges from unresolved internal
structure, including juxtaposed sand and clay units, stochastic clay-smear
continuity, burial and faulting history, and directional permeability.
Representing the fault by one deterministic property field can therefore
understate both the range and the structured, multimodal character of
uncertainty relevant to reservoir-scale forecasts. PREDICT provides a
physics-informed stochastic representation of fault-core architecture and
effective permeability for normally consolidated sediments
\cite{SaloSalgado2025}. Previous work has also demonstrated the importance of
propagating fault-property uncertainty into predictions of \coTwo{} migration
and fault response \cite{Lu2025}.

The remaining challenge is computational scale. An end-to-end assessment must
span plausible geologies, stochastic fault realizations within each throw
window, uncertain relationships among windows, spatial variability along
strike, upscaled multiphase properties, and long-term field-scale simulation.
Direct enumeration of this combined space is infeasible, whereas aggressive
averaging can remove permeability-tensor dependence, distinct probability
modes, and capillary barriers that control migration.

We develop an end-to-end uncertainty quantification workflow with three
computational stages: stochastic fault-property prediction, hierarchical
sampling of three-dimensional fault properties, and high-performance
multiphase reservoir simulation. The sampling hierarchy separates geologic
uncertainty, within-window PREDICT stochasticity, and cross-window coupling
uncertainty. PREDICT produces window-specific distributions but does not
identify a true cross-window joint law; consequently, distributional
similarity is used to construct transparent candidate coupling assumptions,
not to claim measured cross-window correlation.

We apply the workflow to a megatonne-scale offshore Texas storage case with a
three-dimensional growth fault that partially offsets the top seal. The study
uses 162 geologic scenarios and 10 sampled fault-property cases per geology,
yielding 1,620 reservoir simulations over a 1,000-year horizon. We quantify
storage-reservoir retention, entry into the fault and top seal, maximum upward
migration, and partitioning between dissolved and free \coTwo{}. The objectives
are to determine how uncertainty propagates from geologic characterization to
these containment outcomes, identify the dominant controls, and evaluate the
implications for storage decisions. \todo{Add a concise statement of the main
quantitative findings after the production ensemble is complete.}
